\chapter{Conclusion}
\label{ch:conclusion}

This dissertation presents a machine-checked EasyCrypt proof surface for the
provable security of the formal HAETAE random-oracle signature model.  The word
``formal'' is essential: the checked theorem surface is a conditional
implication about the EasyCrypt model, together with an audited correspondence
to the February 2026 HAETAE specification.  The proof is organized as a
game-based reduction with explicit oracle state, transcript state, sampler
state, bad events, query budgets, and algebraic loss composition.

The main contribution is the conversion of informal proof phrases into checked
mathematical objects.  Games are probabilistic programs.  Random oracles are
stateful lazy tables.  Hybrid transitions are equivalence judgments or
probability inequalities.  Bad events are Boolean predicates over concrete game
state.  Active signing has local clean-conditioned one-call proof surfaces and
wrapper invariants; the final NMA-level lift still has to avoid coarse fallback
bounds before it can be read as a non-vacuous active-signing theorem.

The revised dissertation also separates the checked implication from its
boundaries.  The result is conditional on the hardness terms represented in
\ec{HAETAE\_Assumptions.ec}, assembled in \ec{HAETAE\_Reductions.ec}, and on the
explicit hop/reduction antecedents of the final EasyCrypt composition lemma.
One of those antecedents is currently vacuous under the present
\ec{rejection\_sampling\_loss\_term} constants, so the central artifact is a
ROM-model implication rather than a non-vacuous concrete security theorem.  It
also does not prove implementation correctness of all concrete parsing,
arithmetic, hashing, and rejection-sampling code paths.

\section*{Evidence summary}

\begin{longtable}{p{0.28\textwidth}p{0.62\textwidth}}
\toprule
\textbf{Evidence item} & \textbf{Status recorded in this dissertation} \\
\midrule
Specification source & \ec{HAETAE\_v260204.pdf}, 46 pages, generated February
4, 2026; Section 4.2, Figure 8, and Table 1 are used for model/spec
correspondence. \\
\midrule
Proof manifest & \ec{provable-security/proof-files.txt}; recorded inventory:
16 EasyCrypt files and 22046 checked-source lines. \\
\midrule
Main checked bound & \ec{haetae\_euf\_counted\_rom\_bound}; the right-hand side
of the conditional probability-only ROM implication for \(\mathsf{ECModel}\),
with the proof and bound ledger instantiated at \(q_H=q_S=2\), expanded as MLWE hardness, Module-SIS hardness,
bimodal-to-MSIS reduction loss, and counted ROM collision/prequery/min-entropy
terms.
\\
\midrule
Strongest exact claim & A checked implication whose conclusion is a probability
bound for \(\mathsf{ECModel}\) in the ROM; the checked bound ledger uses
\(q_H=q_S=2\).  The bridge from \(\mathsf{Spec}_{260204}\) to
\(\mathsf{ECModel}\) is documented by audit tables and by the explicit
signing-randomness modeling assumption in
\Cref{assm:signing-randomness-modeling-bridge}, not by a separate checked
equivalence theorem.  The current final-hop premise contains
\ec{rejection\_sampling\_loss\_term}, which is larger than one under the present
constants, so this is a proof-surface implication and not a non-vacuous concrete
instantiation. \\
\midrule
Concrete counted-ROM loss & With the current constants in \ec{HAETAE\_Reductions.ec},
\(q_H=2\), \(q_S=2\), \(N=2^{58}\), and the counted ROM contribution is
\(37/2^{58}\). \\
\midrule
Local non-vacuity evidence & Concrete lazy-ROM sampler freshness, sampler-ROM
coverage, \ec{sampler\_bad\_prequery} instrumentation, and active \ec{O.sign}
clean lifting lemmas are identified by name and interpreted mathematically.
The final NMA-level path still contains coarse fallback bounds that must be
replaced before the top-level antecedent can be called non-vacuous. \\
\midrule
Reproducibility aid & \ec{dissertation/etc/count-easycrypt-declarations.sh}
records the counting policy for the proof-code inventory. \\
\midrule
Residual obligations & Replacing the vacuous final-hop antecedent with a
satisfiable checked hop, mechanizing the model/spec fidelity ledger in
\Cref{tab:internal-algorithm-fidelity-ledger}, concrete SHAKE instantiation,
implementation correctness, concrete MLWE/Module-SIS bit-security estimation,
and reduction runtime overhead remain outside the central checked probability
implication. \\
\bottomrule
\end{longtable}

The final lesson is methodological.  A machine-checked cryptographic proof is
not persuasive merely because it is long or because it compiles.  It is
persuasive when the theorem statement is exact, the loss ledger is not inflated,
vacuous hypotheses are called out, the assumptions are visible, and the model is
compared against the specification.  The rewritten dissertation is structured
around those criteria.

The dissertation should therefore be cited in three layers.  First, as a checked
EasyCrypt implication and audit of a scoped counted-ROM proof surface for
\(\mathsf{ECModel}\), with the current vacuity warning attached.  Second, as an
audited specification-model correspondence argument for reading that model as
HAETAE v260204.  Third, as a roadmap for the remaining work needed for an
end-to-end claim: a satisfiable final-hop theorem, mechanized PDF-to-model
equivalence, concrete SHAKE instantiation reasoning, implementation correctness,
concrete MLWE/Module-SIS bit estimates, and cost-aware reduction accounting.
Collapsing these layers would overstate the verification; keeping them separate
is the main integrity property of the document.

Future work should connect the formal HAETAE model to implementation-level
correctness proofs, especially for parsing, hashing, arithmetic, and rejection
sampling.  It should also formalize reduction runtime overhead, maintain a
versioned proof-manifest discipline, and update the model/spec correspondence
whenever the EasyCrypt proof surface or HAETAE specification changes.
